\section{Introduction}
\label{sec:intro}

The proposed Districting Integrity Act (\S~107(c)(1)) mandates that every
state draw its congressional districts by executing a deterministic
recursive-bisection algorithm with a single federally published random seed.
The seed is the statute's reproducibility guarantee: anyone who runs the
reference implementation with the published seed and public census inputs
recovers the identical map, bit for bit. This guarantee is legally and
politically essential. But it raises an immediate challenge:

\begin{quote}
\emph{Why this seed? Did the Director of Census choose it after knowing which
partisan outcome it would produce?}
\end{quote}

This paper answers that challenge rigorously. Our central finding is that
for states where convergence can be certified, the seed choice is irrelevant:

\begin{theorem}[informal]
For each state $s$, define convergence as the event that the running minimum
edge-cut has not improved for $T$ consecutive seeds (empirically: $T \geq
0.5 \times N$ for 1{,}000-seed sweeps). For states where convergence is
certified, the MEC plan is a property of the graph, not of the seed.
Any agent who runs enough seeds finds the same plan.
\end{theorem}

\paragraph{Caveat on convergence.}
Convergence is not guaranteed in a fixed seed budget. Our 1{,}000-seed sweeps
show that some states with initially plausible MEC results at 200 seeds are
overturned by deeper sweeps: Georgia's running minimum dropped 52\,km between
seed 200 (2{,}598\,km) and seed 489 (2{,}546\,km), changing the interim result
from 6D/8R to 7D/7R. The correct protocol is not ``run $N$ seeds'' but
``run seeds until a certified convergence criterion is met.''

\paragraph{Why edge-cut is the correct criterion.}
Minimum edge-cut has no interpretive degrees of freedom. The graph is the
state's TIGER-derived tract adjacency graph, a federal public record. The
edge weights are shared boundary lengths in metres, computed mechanically.
Minimising edge-cut minimises the total length of district boundaries ---
the most direct geometric formalisation of ``compact districts with minimal
artificial boundaries.''

Compactness metrics --- Polsby-Popper, Reock, convex-hull ratio --- are
related but distinct, and they introduce choices that do not arise in
edge-cut minimisation. Which area measure? Land only (ALAND) or
land-plus-water (ALAND + AWATER)? Which shape metric? These choices could,
in principle, be argued to have been selected post-hoc to favour a particular
partisan outcome. Edge-cut has no such attack surface.

Polsby-Popper nevertheless plays an important role in our analysis as a
\emph{post-hoc audit}: after the minimum-edge-cut plan is identified, we
verify it is not pathologically non-compact. The audit is a check, not a
selection criterion. ``We minimised edge-cut and verified the result is
reasonably compact'' is a stronger and more legally durable claim than
``we selected by a weighted combination of edge-cut and compactness.''

\paragraph{Empirical contributions.}
We study Pennsylvania in depth (17 districts, $\sim$4{,}000 census tracts)
across 1{,}100 random seeds and extend the analysis to six additional
states. Our findings:

\begin{enumerate}
  \item \textbf{Dense solution space.} Every seed produces a distinct
        partition --- 1{,}100 seeds, 1{,}100 unique plans (SHA-256
        verified). The solution space is not a small set of canonical
        boundaries.
  \item \textbf{Edge-cut convergence.} The running minimum edge-cut
        improved 12 times across 1{,}100 seeds, last at seed 181.
        The minimum was stable for 919 consecutive seeds afterward.
        The empirical convergence rate supports a formal certificate
        that the global minimum is within $\sim$1\% of the observed best.
  \item \textbf{Partisan stability.} The minimum-edge-cut plan for
        Pennsylvania produces 8 Democratic seats out of 17 (50.6\%
        statewide Dem vote, gap = $-3.5$\,pp). This outcome was reached
        by 6 of the 12 improvement steps. It is not an outlier in the
        partisan distribution: 35\% of all 1{,}100 seeds produce
        the same outcome class (8D).
  \item \textbf{Cross-state validation.} [Results from NC, MI, WI, GA,
        AZ, MA, OH pending; see \S\ref{sec:eval:national}.]
\end{enumerate}

\paragraph{The \textsc{CompactBisect} algorithm.}
A complementary contribution is \textsc{CompactBisect}: a modification that
applies a compactness criterion \emph{at every bisection level} rather than
only at the final partition. At each level, $N$ METIS candidates are
generated; the one maximising $\sqrt{\mathrm{PP}(L) \cdot \mathrm{PP}(R)}$
(geometric-mean Polsby-Popper of both halves) is selected. The geometric
mean is used because it is maximised when both halves are equally compact ---
the same fairness principle that underlies the Huntington-Hill equal-proportions
apportionment method (2~U.S.C.~\S~2a). \textsc{CompactBisect} provides
level-by-level geometric justification for every intermediate decision in
the recursion tree.

Whether \textsc{CompactBisect} produces the same plan as minimum-edge-cut
selection is an open empirical question, and its answer is a central result
of this paper.

\paragraph{Outline.}
\S\ref{sec:background} covers background on graph bisection hardness and
compactness metrics. \S\ref{sec:methodology} defines the minimum-edge-cut
protocol, the \textsc{CompactBisect} algorithm, and our convergence
framework. \S\ref{sec:evaluation} presents empirical results.
\S\ref{sec:discussion} draws implications for the Districting Integrity Act.
\S\ref{sec:conclusion} concludes.
